% ---- score2 standard preamble (from: python3 score2_tools.py preamble) ----
\documentclass[10pt]{extarticle}
\usepackage[a4paper,margin=1.5cm]{geometry}
\usepackage[T1]{fontenc}
\usepackage{amsmath,amssymb}
\usepackage{graphicx}
\usepackage[export]{adjustbox}
\usepackage{multicol}
\usepackage{enumitem}
\usepackage{booktabs,array,tabularx}
\setlength{\columnsep}{0.7cm}
\setlength{\parindent}{0pt}
\setlength{\parskip}{3pt}
\pagestyle{empty}
\makeatletter
% figure registry (figdefs.tex is generated by the extract step)
\newcommand{\qfigdef}[4]{\@namedef{qf@file@#1}{#2}\@namedef{qf@w@#1}{#3}\@namedef{qf@wide@#1}{#4}}
\newcommand{\qf@check}[1]{\@ifundefined{qf@file@#1}{\PackageError{score2}{Unknown figure id #1}{}}{}}
\newcommand{\qf@img}[4][]{\includegraphics[width=#2\dimexpr\@nameuse{qf@w@#4}pt\relax,max width=#3,max height=0.45\textheight,#1]{\@nameuse{qf@file@#4}}}
% Full-width questions: finalize wraps every question that contains a wide figure in \qfullbegin...\qfullend,
% so the WHOLE question (stem, figures, options) is set across the page in reading order, then the two columns
% resume. The question's head - its number, stem and first figure (\qfullsplit marks the end of the head) -
% is kept on one page, so a number is never left at a page foot with its figure overleaf; the rest may break
% across pages normally, so no big gaps are left. Without a \qfullsplit, the whole question is kept together
% if it is under 45% of a page. Never write these yourself.
\newif\ifqf@full
\newif\ifqf@boxed
\newsavebox{\qf@box}
\newcommand{\qf@open}{\global\qf@fulltrue\par\global\qf@boxedtrue
  \setbox\qf@box\vbox\bgroup\hsize\textwidth\linewidth\textwidth}
\newcommand{\qf@emit}[1]{\par\egroup\global\qf@boxedfalse
  \ifdim\dimexpr\ht\qf@box+\dp\qf@box\relax<#1\textheight\noindent\usebox{\qf@box}\par\else\unvbox\qf@box\fi}
\newcommand{\qfullsplit}{\ifqf@boxed\qf@emit{0.85}\fi}
\newcommand{\qf@close}{\ifqf@boxed\qf@emit{0.45}\fi\global\qf@fullfalse}
\newcommand{\qfullbegin}{\end{multicols}\qf@open}
\newcommand{\qfullend}{\qf@close\begin{multicols}{2}}
\newcommand{\qfullstart}{\qf@open}
\newcommand{\qfulldone}{\qf@close}
% \qfig[scale]{id}: block figure, centred, directly where it is written (size set by the tool).
% Fallback if a wide figure is not inside a full-width question: break out of the columns in place - no floats.
\newcommand{\qfig}[2][1]{\qf@check{#2}%
  \ifqf@full\qf@blk{#1}{#2}\else
  \ifnum\@nameuse{qf@wide@#2}=1 \expandafter\@firstoftwo\else\expandafter\@secondoftwo\fi
  {\end{multicols}\begin{center}\qf@img{#1}{\textwidth}{#2}\end{center}\begin{multicols}{2}}%
  {\qf@blk{#1}{#2}}\fi}
\newcommand{\qf@blk}[2]{\par\smallskip{\centering\qf@img{#1}{\linewidth}{#2}\par}\smallskip}
% \qfigi[scale]{id}: inline figure (inside option lists / tables / minipages); never breaks columns.
\newcommand{\qfigi}[2][1]{\qf@check{#2}\qf@img[valign=c]{#1}{\linewidth}{#2}}
\makeatother
% structure macros - lint reads these to check completeness and figure placement
\newcommand{\papersec}[1]{\par\bigskip{\centering\Large\bfseries PAPER-#1\par}\medskip}
\newcommand{\partsec}[2]{\par\medskip\noindent{\large\bfseries PART-#1: #2}\par\smallskip}
\newcommand{\sect}[1]{\par\smallskip\noindent\textbf{#1}\par}
\newcommand{\qn}[1]{\par\smallskip\noindent\textbf{#1.}~}
\newcommand{\paraq}[2]{\par\smallskip\noindent\textbf{Paragraph for Questions #1 and #2}\par}
\newenvironment{qoptions}{\begin{itemize}[leftmargin=1.8em,itemsep=1pt,topsep=2pt,label={}]}{\end{itemize}}
\newcommand{\opt}[2]{\item[(#1)] #2}
% figure registry (generated by score2_tools.py finalize)
\qfigdef{p1_phy_s1_q03_1}{Score-II - Test 03 (05 May 2021)-img/p1_phy_s1_q03_1.png}{101.9}{0}
\qfigdef{p1_phy_s1_q05_1}{Score-II - Test 03 (05 May 2021)-img/p1_phy_s1_q05_1.png}{141.8}{0}
\qfigdef{p1_phy_s1_q06_2}{Score-II - Test 03 (05 May 2021)-img/p1_phy_s1_q06_2.png}{51.1}{0}
\qfigdef{p1_phy_s1_q06_1}{Score-II - Test 03 (05 May 2021)-img/p1_phy_s1_q06_1.png}{92.1}{0}
\qfigdef{p1_phy_s1_q07_1}{Score-II - Test 03 (05 May 2021)-img/p1_phy_s1_q07_1.png}{121.3}{0}
\qfigdef{p1_phy_s1_q08_1}{Score-II - Test 03 (05 May 2021)-img/p1_phy_s1_q08_1.png}{82.9}{0}
\qfigdef{p1_phy_s1_q08_2}{Score-II - Test 03 (05 May 2021)-img/p1_phy_s1_q08_2.png}{82.9}{0}
\qfigdef{p1_phy_s1_q08_3}{Score-II - Test 03 (05 May 2021)-img/p1_phy_s1_q08_3.png}{79.5}{0}
\qfigdef{p1_phy_s1_q08_D}{Score-II - Test 03 (05 May 2021)-img/p1_phy_s1_q08_D.png}{81.2}{0}
\qfigdef{p1_phy_s1_q10_1}{Score-II - Test 03 (05 May 2021)-img/p1_phy_s1_q10_1.png}{121.7}{0}
\qfigdef{p1_phy_s1_q11_1}{Score-II - Test 03 (05 May 2021)-img/p1_phy_s1_q11_1.png}{201.6}{0}
\qfigdef{p1_phy_s1_q13_1}{Score-II - Test 03 (05 May 2021)-img/p1_phy_s1_q13_1.png}{121.7}{0}
\qfigdef{p1_phy_s1_q18_1}{Score-II - Test 03 (05 May 2021)-img/p1_phy_s1_q18_1.png}{89.6}{0}
\qfigdef{p1_phy_s1_q19_1}{Score-II - Test 03 (05 May 2021)-img/p1_phy_s1_q19_1.png}{49.5}{0}
\qfigdef{p1_phy_s1_q20_1}{Score-II - Test 03 (05 May 2021)-img/p1_phy_s1_q20_1.png}{57.8}{0}
\qfigdef{p1_phy_s2_q01_1}{Score-II - Test 03 (05 May 2021)-img/p1_phy_s2_q01_1.png}{113.9}{0}
\qfigdef{p1_phy_s2_q02_1}{Score-II - Test 03 (05 May 2021)-img/p1_phy_s2_q02_1.png}{153.6}{0}
\qfigdef{p1_phy_s2_q08_1}{Score-II - Test 03 (05 May 2021)-img/p1_phy_s2_q08_1.png}{95.1}{0}
\qfigdef{p1_phy_s2_q10_1}{Score-II - Test 03 (05 May 2021)-img/p1_phy_s2_q10_1.png}{95.5}{0}
\qfigdef{p1_chem_s1_q04_1}{Score-II - Test 03 (05 May 2021)-img/p1_chem_s1_q04_1.png}{181.7}{0}
\qfigdef{p1_chem_s1_q12_2}{Score-II - Test 03 (05 May 2021)-img/p1_chem_s1_q12_2.png}{58.3}{0}
\qfigdef{p1_chem_s1_q12_1}{Score-II - Test 03 (05 May 2021)-img/p1_chem_s1_q12_1.png}{59.8}{0}
\qfigdef{p1_chem_s1_q12_3}{Score-II - Test 03 (05 May 2021)-img/p1_chem_s1_q12_3.png}{33.2}{0}
\qfigdef{p1_chem_s2_q03_1}{Score-II - Test 03 (05 May 2021)-img/p1_chem_s2_q03_1.png}{102.4}{0}
\qfigdef{p1_chem_s2_q04_1}{Score-II - Test 03 (05 May 2021)-img/p1_chem_s2_q04_1.png}{79.8}{0}
\qfigdef{p1_chem_s2_q04_2}{Score-II - Test 03 (05 May 2021)-img/p1_chem_s2_q04_2.png}{31.2}{0}
\qfigdef{p1_chem_s2_q04_3}{Score-II - Test 03 (05 May 2021)-img/p1_chem_s2_q04_3.png}{38.8}{0}
\qfigdef{p1_chem_s2_q04_4}{Score-II - Test 03 (05 May 2021)-img/p1_chem_s2_q04_4.png}{31.2}{0}
\qfigdef{p1_chem_s2_q04_5}{Score-II - Test 03 (05 May 2021)-img/p1_chem_s2_q04_5.png}{147.7}{0}
\raggedcolumns
\makeatletter\@beginparpenalty=10000\makeatother
\begin{document}
\begin{center}{\Large\bfseries JEE (Main+Advanced) Score-II --- Test 03 (05 May 2021)}\end{center}
\begin{multicols}{2}

\papersec{1}
\partsec{1}{PHYSICS}
\sect{Section-I}

\qn{1} Determine the change in entropy $\Delta S$ of 100 gm of water when it is heated to the boiling point from room temperature of T = 300 K, and it is transformed into vapor at the boiling temperature and atmospheric pressure. It is given that for water : the molar heat capacity C = 75.4 J/(mol-K); latent heat of vaporization L = 2.26 MJ/kg.\\
\relax[Take ln(1.24) = 0.217]
\begin{qoptions}
\opt{A}{697 J/K}
\opt{B}{600 J/K}
\opt{C}{500 J/K}
\opt{D}{650 J/K}
\end{qoptions}

\qn{2} Oscillator is moving according to the law, $x=x_0\sin\omega t$ because of a driving force $F=F_0\cos\omega t$ acting on it. What is the damping coefficient $\gamma$ of the oscillator ? Mass of the oscillator is m. The resistive force is dependent on velocity v as $F=-\gamma mv$.
\begin{qoptions}
\opt{A}{$\dfrac{F_o}{x_o\omega}$}
\opt{B}{$\dfrac{F_o}{2x_o\omega}$}
\opt{C}{$\dfrac{F_o}{mx_o\omega}$}
\opt{D}{$\dfrac{F_o}{2mx_o\omega}$}
\end{qoptions}

\qn{3} A small particle of mass m is at rest on a horizontal circular platform that is free to rotate about a vertical axis through its center. The particle is located at a radius r from the axis, as shown in the figure. The platform begins to rotate with constant angular acceleration $\alpha$. Because of friction between the particle and the platform, the particle remains at rest with respect to the platform. When the platform has reached angular speed $\omega$, the angle $\theta$ between the static frictional force $f_s$ and the inward radial direction is :-
\qfig{p1_phy_s1_q03_1}
\begin{qoptions}
\opt{A}{$\theta=\dfrac{\omega^2r}g$}
\opt{B}{$\theta=\dfrac\alpha{\omega^2}$}
\opt{C}{$\theta=\tan^{-1}\left(\dfrac{\omega^2}\alpha\right)$}
\opt{D}{$\theta=\tan^{-1}\left(\dfrac\alpha{\omega^2}\right)$}
\end{qoptions}

\qn{4} Which of the following sample of gases has the same internal energy as the energy of one mole of hydrogen gas at 27$^\circ$C :-
\begin{qoptions}
\opt{A}{1 mole of Helium gas at 327$^\circ$C}
\opt{B}{3 mole of Helium gas at 27$^\circ$C}
\opt{C}{0.5 mole of O$_2$ gas at 327$^\circ$C}
\opt{D}{2 mole of O$_2$ gas at 27$^\circ$C}
\end{qoptions}

\qn{5} Thin tape of width $\ell=40$ cm is rolled into a tube of radius R = 30 cm. On the tape flows evenly distributed current of I = 200 A as shown in figure. The magnetic induction B on the axis of the tube at point 2 is :-
\qfig{p1_phy_s1_q05_1}
\begin{qoptions}
\opt{A}{100 $\pi\mu$T}
\opt{B}{80 $\pi\mu$T}
\opt{C}{60 $\pi\mu$T}
\opt{D}{120 $\pi\mu$T}
\end{qoptions}

\qn{6} If gravitational force between two identical balls is F as shown in figure I. What is the force between two balls when arranged as in figure II.
\begin{center}\qfigi{p1_phy_s1_q06_2}\qquad\qfigi{p1_phy_s1_q06_1}\end{center}
\begin{qoptions}
\opt{A}{$\left(1+\dfrac{\sqrt3}2\right)$F}
\opt{B}{F}
\opt{C}{$\dfrac32$F}
\opt{D}{2F}
\end{qoptions}

\qn{7} Two-small magnetic dipole with the same magnetic moments $\mu$ are located at a distance r from each other. Determine force of their interaction.
\qfig{p1_phy_s1_q07_1}
\begin{qoptions}
\opt{A}{$F=\dfrac{3\mu_o\mu^2}{2\pi\,r^4}$}
\opt{B}{$F=\dfrac{\mu_o\mu^2}{\pi\,r^4}$}
\opt{C}{$F=\dfrac{1\mu_o\mu^2}{2\pi\,r^4}$}
\opt{D}{$F=\dfrac{\mu_o\mu^2}{2\pi\,r^4}$}
\end{qoptions}

\qn{8} A point mass m is connected to an ideal spring on a horizontal frictionless surface. The mass is pulled a short distance and then released. Which of the following is the most correct plot of the kinetic energy as a function of potential energy ?
\begin{qoptions}
\opt{A}{\qfigi{p1_phy_s1_q08_1}}
\opt{B}{\qfigi{p1_phy_s1_q08_2}}
\opt{C}{\qfigi{p1_phy_s1_q08_3}}
\opt{D}{\qfigi{p1_phy_s1_q08_D}}
\end{qoptions}

\qn{9} The potential energy function between two atoms in a diatomic molecule can be expressed approximately as follows $U(x)=\dfrac a{x^{12}}-\dfrac b{x^6}$, where a and b are positive constants and x is the distance between atoms. What is the dissociation energy of the molecule?
\begin{qoptions}
\opt{A}{$\dfrac{b^2}a$}
\opt{B}{$\dfrac{b^2}{2a}$}
\opt{C}{$\dfrac{b^2}{4a}$}
\opt{D}{$\dfrac{2b^2}a$}
\end{qoptions}

\qn{10} Isosceles glass prism ($\mu$ = 1.57) with small refracting angles ($\alpha$) is placed in front of parallel beam of rays incident normally on its base (see figure). The size of the base 2a = 5 cm. Light falls on a screen at a distance of 100 cm from the prism with a dark patch of width 1 cm near center of screen. Find the value of the refracting angle $\alpha$.
\qfig{p1_phy_s1_q10_1}
\begin{qoptions}
\opt{A}{$3^\circ$}
\opt{B}{$4^\circ$}
\opt{C}{$5^\circ$}
\opt{D}{$6^\circ$}
\end{qoptions}

\qn{11} A uniform cylinder of radius a originally has a weight of 80 N. After an off-axis cylindrical hole, at a distance of 2a/5 from centre was drilled through it, it weighs 65 N. The axes of the two cylinders are parallel and their centers are at the same height. A force T is applied to the top of the cylinder horizontally. In order to keep the cylinder at rest, the magnitude of the force is closest to:
\qfig{p1_phy_s1_q11_1}
\begin{qoptions}
\opt{A}{3 N}
\opt{B}{6N}
\opt{C}{15 N}
\opt{D}{30 N}
\end{qoptions}

\qn{12} The correct relation between the angle of diffraction $\phi$ and the glancing angle $\theta$ in Davission--Germer experiment will be:
\begin{qoptions}
\opt{A}{$\theta=90^\circ-\dfrac\phi2$}
\opt{B}{$\phi=\dfrac\theta2-90^\circ$}
\opt{C}{$\theta=90^\circ-\phi$}
\opt{D}{None of these}
\end{qoptions}

\qn{13} Both the blocks shown here are of mass m and are moving with constant velocity in direction shown in a resistive medium which exerts equal constant force on both blocks in direction opposite to the velocity. The tension in the string connecting both of them will be: (Neglect friction)
\qfig{p1_phy_s1_q13_1}
\begin{qoptions}
\opt{A}{mg}
\opt{B}{mg/2}
\opt{C}{mg/3}
\opt{D}{mg/4}
\end{qoptions}

\qn{14} A projectile is fired with a speed of $\sqrt{\dfrac{GM}r}$ at certain angle ($\theta\neq0^\circ,90^\circ,180^\circ$) with the radius vector. Here r is distance of projectile from center of Earth and M is mass of Earth. Choose the \textbf{INCORRECT} statement.
\begin{qoptions}
\opt{A}{The projectile will follow elliptical path.}
\opt{B}{The projectile is at semi-minor axis of its elliptical path.}
\opt{C}{Earth is at nearer focus if the projectile is initially approaching the Earth.}
\opt{D}{Earth could be at any focus of the ellipse.}
\end{qoptions}

\qn{15} Consider Fraunhofer diffraction pattern obtained with a single slit illuminated at normal incidence. At the angular position of the first diffraction minimum, the phase difference (in radians) between the wavelets from the opposite edges of the slit is :-
\begin{qoptions}
\opt{A}{$\pi$}
\opt{B}{$2\pi$}
\opt{C}{$\pi/4$}
\opt{D}{$\pi/2$}
\end{qoptions}

\qn{16} A boat is traveling due east at 12 m/s. A flag on the boat flaps at 53$^\circ$N of W. Another flag on the shore flaps due north :-
\begin{qoptions}
\opt{A}{Speed of wind with respect to ground is 16 m/s}
\opt{B}{Speed of wind with respect to ground is 20 m/s}
\opt{C}{Speed of wind with respect to boat is 24 m/s}
\opt{D}{Speed of wind with respect to boat is 16 m/s}
\end{qoptions}

\qn{17} A paper helicopter with rotor radius r and weight W is dropped from a height h in air with a density of $\rho$. Assuming that the helicopter quickly reaches terminal velocity, a function for the time of flight T can be found in the form\\
$T=kh^\alpha r^\beta\rho^\delta W^\omega$\\
where k is an unknown dimensionless constant. $\alpha$, $\beta$, $\delta$ and $\omega$ are constant exponents to be determined. Determine $\alpha$.
\begin{qoptions}
\opt{A}{$\alpha=-1$}
\opt{B}{$\alpha=-1/2$}
\opt{C}{$\alpha=0$}
\opt{D}{$\alpha=1$}
\end{qoptions}

\qn{18} Consider a system of n concentric conducting shells of radii r, 2r, 4r, 8r, ....., $2^{n-1}$r respectively. If inner shell is earthed and Q charge is given on outer most shell, charge appears on inner most shell is :-
\qfig{p1_phy_s1_q18_1}
\begin{qoptions}
\opt{A}{$-$Q}
\opt{B}{$-\dfrac Q{2^n}$}
\opt{C}{$-\dfrac Q{2^{n-1}}$}
\opt{D}{$+$Q}
\end{qoptions}

\qn{19} Figure shows a 5 kg ladder hanging from a string that is connected with a ceiling and is having a spring balance connected in between. A man of mass 25 kg is climbing up the ladder at acceleration 1 m/s$^2$. Assuming the spring balance and the string to be massless, the reading of the spring balance is :-
\qfig{p1_phy_s1_q19_1}
\begin{qoptions}
\opt{A}{30 kg}
\opt{B}{32.5 kg}
\opt{C}{35 kg}
\opt{D}{37.5 kg}
\end{qoptions}

\qn{20} On a smooth level ground we keep a light rod to which 2 masses M \& 2M are attached. The velocities of these mass at the moment is shown here. Choose the \textbf{INCORRECT} statement.
\qfig{p1_phy_s1_q20_1}
\begin{qoptions}
\opt{A}{The velocities of masses remain constant.}
\opt{B}{The angular velocity of the rod is V/L clockwise}
\opt{C}{The rod is under tension}
\opt{D}{The centre of mass will move in a straight line.}
\end{qoptions}

\sect{Section-II}

\qn{1} A hemispherical bowl of inner radii R and outer radii 2R is charged with density $\rho=\rho_o\dfrac rR$. The electric field at point O is $\dfrac{b\rho_oR}{\varepsilon_o}$. The value of b is :
\qfig{p1_phy_s2_q01_1}

\qn{2} The current (in Ampere) through ammeter A in circuit shown is :
\qfig{p1_phy_s2_q02_1}

\qn{3} A distant hot air balloon subtends $0.25^\circ$ at the objective lens of an astronomical telescope. The image subtends $1.5^\circ$. The objective and eye piece are in a 35 cm long tube. The focal length (in m) of the eye piece lens is

\qn{4} The period of oscillation of a magnet in a vibration magnetometer is 2 seconds. The period of oscillate (in sec) of a magnet of moment of inertia four times that of the first magnet is :-

\qn{5} A ball is thrown vertically downwards with velocity $\sqrt{2gh}$ from a height h. After colliding with the ground it just reaches the starting point. Coefficient of restitution is e. The value of $e^2$ is

\qn{6} A Fresnel biprism is used to form the interference fringes. The distance between the source and the biprism is 20 cm and that between the biprism and the screen is 80 cm. If $\lambda$ = 6563 \AA{} and the separation between the virtual sources is 3.6 mm, then the fringe width (in mm) is.

\qn{7} Determine the change in entropy $\Delta S$ (in SI unit) of 100 gm of water when it is heated to the boiling point from room temperature of T = 300 K, and it is transformed into vapour at the boiling temperature and atmospheric pressure. It is given that for water : the molar heat capacity C = 75.4 J/(mol-K); latent heat of vaporization: L = 2.26 MJ/kg.\\
\relax[Take $\ell$n(1.24) = 0.217]

\qn{8} In the figure shown, the minimum force F (in N) to be applied perpendicular to the incline so that the block does not slide is
\qfig{p1_phy_s2_q08_1}

\qn{9} Electric potential in space\\
$V(x,y,z)=\left(x^2+y^2+z^2\right)$V\\
Then the magnitude of electric field strength E at (1 m, 0, 2m) in N/C is :

\qn{10} A jet travelling at a constant speed of $1.20\times10^2$ m/s executes a vertical loop with a radius of $5.00\times10^2$ m. (See Fig.) Find the magnitude of the force (in N) of the seat on a 70.0 kg pilot at the top of the loop.\\
(Take : g = 10 m/s$^2$)
\qfig{p1_phy_s2_q10_1}

\partsec{2}{CHEMISTRY}
\sect{Section-I}

\qn{1} In which of the following set of orbitals, no. of spherical nodes are same.
\begin{qoptions}
\opt{A}{$4d_{z^2}$ \& 4f}
\opt{B}{$3d_{x^2-y^2}$ \& $3p_x$}
\opt{C}{$4d_{xy}$ \& $3d_{z^2}$}
\opt{D}{$3p_x$ \& $4d_{x^2-y^2}$}
\end{qoptions}

\qn{2} If specific volume of H$_2$ gas at 40 atm (high pressure) is 2 lit/g, then temperature of the gas assuming it vander waal real gas having a = 0.034 dm$^6$ atm mol$^{-2}$ \& b = 0.2 lit/mol is \underline{\hspace{1cm}}approx.
\begin{qoptions}
\opt{A}{1951 K}
\opt{B}{1854 K}
\opt{C}{878 K}
\opt{D}{975 K}
\end{qoptions}

\qn{3} Select the reactions in which $K_p>K_c$.\\
I. $2\text{NH}_3\underset{}{\overset{T=400K}{\rightleftharpoons}}\text{N}_2(g)+3\text{H}_2(g)$\\
II. $2\text{A}(g)\overset{T=10K}{\rightleftharpoons}2\text{B}(g)+\text{C}(g)$\\
III. $\text{SOCl}_2(g)\overset{450K}{\rightleftharpoons}\text{SO}_2(g)+\text{Cl}_2(g)$\\
IV. $\text{NH}_4\text{NO}_2(s)\overset{370K}{\rightleftharpoons}\text{N}_2(g)+2\text{H}_2\text{O}(\ell)$
\begin{qoptions}
\opt{A}{All}
\opt{B}{I, III, IV}
\opt{C}{I, II, III}
\opt{D}{I \& III}
\end{qoptions}

\qn{4} Select the correct statements regarding following process.
\qfig{p1_chem_s1_q04_1}
I. Solubility of Zn(OH)$_2$ is $10^{-4}$ mol/lit.\\
II. On addition of small quantity of OH$^-$ solubility decreases\\
III. On addition of conc. HCl solubility increases\\
IV. In presence of 0.01M ZnCl$_2$, solubility becomes $10^{-5}$ mol/lit.
\begin{qoptions}
\opt{A}{I \& II}
\opt{B}{I, II \& IV}
\opt{C}{I \& IV}
\opt{D}{I, II, III \& IV}
\end{qoptions}

\qn{5} In which of the following given sol or sol made in the process is positively charged.
\begin{qoptions}
\opt{A}{$\text{FeCl}_3+\text{H}_2\text{O}\rightarrow\underset{\text{sol}}{\text{Fe(OH)}_3}+\text{H}^++\text{Cl}^-$}
\opt{B}{$\text{As}_2\text{O}_3+\text{H}_2\text{S}\rightarrow\underset{\text{sol}}{\text{As}_2\text{S}_3}+\text{H}_2\text{O}$}
\opt{C}{$\text{AgNO}_3+\underset{\text{excess}}{\text{KI}}\rightarrow\underset{\text{sol}}{\text{AgI}}+\text{KNO}_3\text{(aq)}$}
\opt{D}{Au sol}
\end{qoptions}

\qn{6} Which of the following surfactant will form micelles in aqueous solution at lowest molar concentration.
\begin{qoptions}
\opt{A}{$\text{CH}_3-\underset{\text{CH}_3}{\text{CH}}-(\text{CH}_2)_{12}-\text{COO}^-\text{Na}^+$}
\opt{B}{$\text{CH}_3-\underset{\text{CH}_3}{\text{CH}}-(\text{CH}_2)_{15}-\text{OSO}_3^-\text{K}^+$}
\opt{C}{$\text{CH}_3-\underset{\text{CH}_3}{\text{CH}}-(\text{CH}_2)_{10}-\text{CH}_2-\text{NH}_3^+\text{Br}^-$}
\opt{D}{$\text{CH}_3-\underset{\text{CH}_3}{\text{CH}}-(\text{CH}_2)_{11}-\text{COO}^-\text{K}^+$}
\end{qoptions}

\qn{7} EMF of the following cell at 25$^\circ$C is\\
Pt $|$ H$^+$ (1M), MnO$_4^-$ (0.1M), Mn$^{+2}$ (0.1M) $||$ H$_2$ (1 atm) $|$ H$_2$SO$_4$ (0.05M)$|$ Pt\\
(given that $E^\circ_{\text{H}^+,\text{Mn}^{+2},\text{MnO}_4^-}=-1.51$V at 25$^\circ$C. 2.303RT/F = 0.06)
\begin{qoptions}
\opt{A}{1.39 V}
\opt{B}{1.57 V}
\opt{C}{1.52 V}
\opt{D}{$-$1.42 V}
\end{qoptions}

\qn{8} Which of the following is correct?
\begin{qoptions}
\opt{A}{D(+)--glucose and D(+)--mannose are enantiomers.}
\opt{B}{D(+)--glucose and D($-$)--fructose give same osazone.}
\opt{C}{D($-$)--fructose give single product on reduction with Na--Hg.}
\opt{D}{D(+) --glucose and D($-$) --fructose are epimers.}
\end{qoptions}

\qn{9} Which one of the following reagent can be used for given conversion?\\
$\text{CH}_3\text{COOCH}_2\text{CH}_3\longrightarrow2\text{CH}_3\text{CH}_2\text{OH}$ \hfill (only organic product)
\begin{qoptions}
\opt{A}{NaBH$_4$/EtOH}
\opt{B}{(i) LiAlH$_4$/ether. (ii) H$_3$O$^+$}
\opt{C}{(i) CH$_3$MgBr \ (ii) H$_3$O$^+$}
\opt{D}{I$_2$/NaOH}
\end{qoptions}

\qn{10} Which of the following reagent/series of reagents can be used for distinction between methyl alcohol and ethyl alcohol?
\begin{qoptions}
\opt{A}{Conc. HCl/ZnCl$_2$}
\opt{B}{CH$_3$COOH/H$^+$}
\opt{C}{NaOH/I$_2$}
\opt{D}{Red P, Br$_2$/AgNO$_2$/HNO$_2$/NaOH}
\end{qoptions}

\qn{11} Which of the following reaction is used to produce aldehyde?
\begin{qoptions}
\opt{A}{$\text{CH}_3\text{CH}_2\text{COOH}\xrightarrow[\text{ii)H}_3\text{O}^+]{\text{i)LiAlH}_4/\text{ether}}$}
\opt{B}{$\text{CH}_3\text{COCl}\xrightarrow{\text{R}_2\text{CuLi}}$}
\opt{C}{$\text{CH}_3\text{COOH (vapours)}\xrightarrow[\Delta]{\text{MnO}}$}
\opt{D}{$\text{CH}_3\text{CN}\xrightarrow[\text{ii)H}_3\text{O}^+]{\text{i)SnCl}_2/\text{HCl}}$}
\end{qoptions}

\qn{12} Which of the following compound can exhibit optical isomerism?
\begin{qoptions}
\opt{A}{\qfigi{p1_chem_s1_q12_2}}
\opt{B}{\qfigi{p1_chem_s1_q12_1}}
\opt{C}{\qfigi{p1_chem_s1_q12_3}}
\opt{D}{CH$_3$--CH=C=CH$_2$}
\end{qoptions}

\qn{13} Which of the following compound cannot be synthesized in the given reaction?\\
$\text{CH}_3-\text{CH}=\text{CH}-\underset{\text{CH}_3}{\text{CH}}-\text{CH}_3\xrightarrow[h\nu]{\text{NBS/CCl}_4}$
\begin{qoptions}
\opt{A}{$\text{CH}_3-\text{CH}=\text{CH}-\underset{\text{CH}_2\text{Br}}{\text{CH}}-\text{CH}_3$}
\opt{B}{$\text{BrCH}_2-\text{CH}=\text{CH}-\underset{\text{CH}_3}{\text{CH}}-\text{CH}_3$}
\opt{C}{$\text{CH}_3-\text{CH}=\text{CH}-\underset{\text{CH}_3}{\overset{\text{Br}}{\text{C}}}-\text{CH}_3$}
\opt{D}{$\text{CH}_3-\underset{\text{Br}}{\text{CH}}-\text{CH}=\underset{\text{CH}_3}{\text{C}}-\text{CH}_3$}
\end{qoptions}

\qn{14} The reaction $[\text{Fe(CN)}_6]^{3-}\rightarrow[\text{FeF}_6]^{3-}$ takes place with :
\begin{qoptions}
\opt{A}{Decrease in magnetic moment}
\opt{B}{Increase in magnetic moment}
\opt{C}{Decrease in co-ordination number}
\opt{D}{Increase in co-ordination number}
\end{qoptions}

\qn{15} Which of the following is not hydrolysed?
\begin{qoptions}
\opt{A}{CCl$_4$}
\opt{B}{SiCl$_4$}
\opt{C}{SnCl$_4$}
\opt{D}{PbCl$_4$}
\end{qoptions}

\qn{16} Which of the following is not planar?
\begin{qoptions}
\opt{A}{H$_2$O}
\opt{B}{BrF$_3$}
\opt{C}{XeF$_4$}
\opt{D}{None of these}
\end{qoptions}

\qn{17} Which one of the following gives metal hydroxide, H$_2$O$_2$ and O$_2$, on hydrolysis, at 0$^\circ$C?
\begin{qoptions}
\opt{A}{BeO}
\opt{B}{Li$_2$O}
\opt{C}{KO$_2$}
\opt{D}{Na$_2$O$_2$}
\end{qoptions}

\qn{18} Which of the following is of red colour?
\begin{qoptions}
\opt{A}{ZnS}
\opt{B}{PbS}
\opt{C}{Ag$_2$CrO$_4$}
\opt{D}{HgO}
\end{qoptions}

\qn{19} The ore that is concentrated by froth floatation process is :
\begin{qoptions}
\opt{A}{Cinnabar}
\opt{B}{Calamine}
\opt{C}{Carnallite}
\opt{D}{Zincite}
\end{qoptions}

\qn{20} Phosphine, ammonia and acetylene can be formed by treating with water :
\begin{qoptions}
\opt{A}{Mg$_3$P$_2$, Li$_3$N, Al$_4$C$_3$}
\opt{B}{Ca$_3$P$_2$, Mg$_3$N$_2$, CaC$_2$}
\opt{C}{Ca$_3$P$_2$, CaCN$_2$, Be$_2$C}
\opt{D}{Ca$_3$P$_2$, NH$_4$NO$_3$, Mg$_2$C$_3$}
\end{qoptions}

\sect{Section-II}

\qn{1} Two mole of an ideal gas is isothermally compressed in two stages starting from ( 1 atm, 20 lit) to 10 lit. at external pressure $P_1$ then upto 2 lit at external pressure $P_2$. Find work done (lit-atm) in the process.

\qn{2} If Osmotic pressure (in atm) of 0.01 M solution at 27$^\circ$C, containing weak acid (HA) having $K_a=6.4\times10^{-5}$ is x. Find the value of (x $\times$ 10) (Given R = 0.082 L-atm/K-mole)

\qn{3} Calculate number of moles of HI required to react with one mole of the given compound -
\qfig{p1_chem_s2_q03_1}

\qn{4} Among the following compound, find the number of compounds which is/are more basic than aniline
\begin{center}\qfigi{p1_chem_s2_q04_1}\quad\qfigi{p1_chem_s2_q04_2}\quad\qfigi{p1_chem_s2_q04_3}\\[2pt]\qfigi{p1_chem_s2_q04_4}\quad\qfigi{p1_chem_s2_q04_5}\end{center}

\qn{5} The brown ring complex compound is formulated as $[\text{Fe(H}_2\text{O)}_5\text{(NO)}]\text{SO}_4$. Find the oxidation state of Fe in the complex.

\qn{6} Glucose + NH$_2-$NH$-$Ph $\longrightarrow$ Product\\
(1 mol) \hspace{1cm} (X gram)\\
Considering reaction took place with 100\% yield of major product find the value of X

\qn{7} For the galvanic cell :\\
Fe(s) $|$ Fe$^{3+}$(0.02M) $||$ Cu$^{2+}$ (0.1M) $|$ Cu (s)\\
the EMF of cell at 298 K is 0.404 V. The value of $\Delta G^o$, (in kJ) for the cell reaction per mole of Fe(s) reacted at 298 K is\\
$\left[\dfrac{2.303\times R\times298}{96500}=0.06,\ 1F=96500\,C\right]$\\
(log2 = 0.3)

\qn{8} Liquids `A' and `B' form an ideal solution. If the mole - fraction of `A' in liquid and vapour forms at equiilibrium is 0.4 and 0.6, respectively, the value of $\left(\dfrac{P_A^0}{P_B^0}\right)$ is.

\qn{9} How many of the following statements is/are \textbf{INCORRECT} regarding ortho hydrogen and para hydrogen?\\
(i) number of neutron in ortho hydrogen is greater than that in para hydrogen\\
(ii) nuclear spins are parallel in para hydrogen while opposite in ortho hydrogen\\
(iii) classification of these two is done based upon electronic spin.\\
(iv) chemical properties are same for both.\\
(v) physical properties are different for them\\
(vi) para hydrogen is more stable compared to ortho hydrogen\\
(vii) conversion of 100\% ortho hydrogen is possible at high temperature

\qn{10} How many of the following is/are the correct match.\\
(i) Photochemical smog : CF$_2$Cl$_2$ (common component)\\
(ii) Ozone layer : Protects from the Sun's radiation of wavelength 255 nm\\
(iii) Methemoglobinemia : Caused by the presence of greater than 50 ppm of nitrate in drinking water\\
(iv) Ozone layer depletion : ClO$^\bullet$(g) + O(g) $\rightarrow$ Cl$^\bullet$(g) + O$_2$(g) (one of the reaction involved)\\
(v) DDT : Biodegradable pollutant\\
(vi) Green house gas : Methane

\partsec{3}{MATHEMATICS}
\sect{Section-I}

\qn{1} If equation of common tangent to curves $y^2=4(x+1)$, $x^2=4(y+1)$ and $x^2+y^2=2$ is $ax+by+4=0$ and A is area of triangle formed by line $ax+by+4=0$ with co-ordinate axes, then
\begin{qoptions}
\opt{A}{A = 1sq.unit}
\opt{B}{A = 2sq.unit}
\opt{C}{A = 3sq.unit}
\opt{D}{A = 4sq.unit}
\end{qoptions}

\qn{2} If $f(x)=x^3+ax^2+b$ has three real and distinct roots and $f'(2)=0$. Then b can't take the value
\begin{qoptions}
\opt{A}{1}
\opt{B}{2}
\opt{C}{3}
\opt{D}{5}
\end{qoptions}

\qn{3} If the normals at the points $P(x_i,y_i)$, i = 1, 2, 3, 4 on the hyperbola xy = 4 are concurrent at the point Q(3, 2) then $\dfrac{(x_1+x_2+x_3+x_4)(y_1+y_2+y_3+y_4)}{x_1x_2x_3x_4}$ is equal to
\begin{qoptions}
\opt{A}{$\dfrac{\sqrt6}{16}$}
\opt{B}{$-\dfrac9{16}$}
\opt{C}{$-\dfrac38$}
\opt{D}{$\dfrac38$}
\end{qoptions}

\qn{4} Let $\dfrac{x^2}{144}+\dfrac{y^2}{81}=1$ be the equation of an ellipse. If m be the minimum value of the segment of a tangent to the ellipse intercepted by the coordinate axes and M be the maximum distance of the centre of the ellipse from the normal, then the value of mM is equal to
\begin{qoptions}
\opt{A}{63}
\opt{B}{225}
\opt{C}{21}
\opt{D}{108}
\end{qoptions}

\qn{5} A line is inclined at angles $\alpha,\beta,\gamma$ with the positive direction of x, y, z axes respectively. If $\beta+\gamma=\dfrac\pi2$, then the maximum possible value of $(\cos\alpha+\sin\beta+\sin\gamma)^2$ is equal to.
\begin{qoptions}
\opt{A}{1}
\opt{B}{2}
\opt{C}{3}
\opt{D}{None of these}
\end{qoptions}

\qn{6} $\begin{vmatrix}3x^2+19&2x-37&2\\x^2-29&x+31&3\\2x^2+17&3x&5\end{vmatrix}=ax^3+bx^2+cx+d$, then value of a is
\begin{qoptions}
\opt{A}{2}
\opt{B}{5}
\opt{C}{$-5$}
\opt{D}{$-8$}
\end{qoptions}

\qn{7} If $3(\sin^{-1}x)+4(\cos^{-1}y)-2(\cos^{-1}z)=\dfrac{11\pi}2$ then find value of $x^{2015}+y^{2017}+z^{2019}$
\begin{qoptions}
\opt{A}{1}
\opt{B}{2}
\opt{C}{3}
\opt{D}{None of these}
\end{qoptions}

\qn{8} If $\theta\in[0,2\pi]$ then number of orderered pairs (x, $\theta$) which satisfies $x^4-2x^2\cos\theta+1=0$ is/are
\begin{qoptions}
\opt{A}{1}
\opt{B}{2}
\opt{C}{3}
\opt{D}{4}
\end{qoptions}

\qn{9} Let $x_1,x_2\dots\dots x_n$ be n observations such that $\Sigma x_i^2=200$ and $\Sigma x_i=40$ then a possible value of n among the following is
\begin{qoptions}
\opt{A}{4}
\opt{B}{5}
\opt{C}{6}
\opt{D}{10}
\end{qoptions}

\qn{10} If in a frequency distribution, the mean and median are 25 and 26, then its mode is approximately.
\begin{qoptions}
\opt{A}{28}
\opt{B}{26}
\opt{C}{24}
\opt{D}{20}
\end{qoptions}

\qn{11} The negation of $p\leftrightarrow q$ is equivalent to
\begin{qoptions}
\opt{A}{$\sim p\leftrightarrow\sim q$}
\opt{B}{$\sim p\leftrightarrow q$}
\opt{C}{$\sim p\wedge\sim q$}
\opt{D}{None of these}
\end{qoptions}

\qn{12} The total number of reflexive relations for a set with n elements is
\begin{qoptions}
\opt{A}{$2^{n^2-n}$}
\opt{B}{$2^n$}
\opt{C}{$n^2$}
\opt{D}{$2^{n^2}$}
\end{qoptions}

\qn{13} If $U_n=\displaystyle\int_0^{\pi/2}x^n\sin x\,dx$ then the value $U_{10}+90U_8$ is
\begin{qoptions}
\opt{A}{$10(\pi/2)^9$}
\opt{B}{$9(\pi/2)^9$}
\opt{C}{$9(\pi/2)^8+(\pi/2)^9$}
\opt{D}{None of these}
\end{qoptions}

\qn{14} If $\displaystyle\int_0^8f(x)dx=3$ then $\displaystyle\sum_{k=1}^8\int_0^1f(k-1+x)dx$ is
\begin{qoptions}
\opt{A}{3}
\opt{B}{24}
\opt{C}{108}
\opt{D}{None of these}
\end{qoptions}

\qn{15} Let f(x) be a continuous function on R. If $\displaystyle\int_0^1[f(x)-f(2x)]dx=5$ and $\displaystyle\int_0^2[f(x)-f(4x)]dx=10$ then the value of $\displaystyle\int_0^1[f(x)-f(8x)]\,dx$ is equal to
\begin{qoptions}
\opt{A}{15}
\opt{B}{10}
\opt{C}{5}
\opt{D}{None of these}
\end{qoptions}

\qn{16} If $f(x)=f(a-x)$, $g(x)=g(a-x)$ and $3h(x)-4h(a-x)=5$.\\
$\displaystyle\int_0^af(x)g(x)h(x)dx=k\int_0^af(x)g(x)dx$ then the value of k is
\begin{qoptions}
\opt{A}{5}
\opt{B}{$-5$}
\opt{C}{4}
\opt{D}{$-4$}
\end{qoptions}

\qn{17} Let $\alpha=\displaystyle\lim_{n\to\infty}\dfrac{1^2+2^2+\dots\dots+n^2}{n^3}$ and\\
$\beta=\displaystyle\lim_{n\to\infty}\dfrac{(1^3-1^2)+(2^3-2^2)+\dots\dots+(n^3-n^2)}{n^4}$ then
\begin{qoptions}
\opt{A}{$\alpha=\beta$}
\opt{B}{$4\alpha=3\beta$}
\opt{C}{$3\alpha=4\beta$}
\opt{D}{None of these}
\end{qoptions}

\qn{18} If $\displaystyle\lim_{x\to0}f(x)=a^3,a\neq0$, then $\displaystyle\lim_{x\to0}f\left(\dfrac xa\right)$ is
\begin{qoptions}
\opt{A}{$a^2$}
\opt{B}{$1/a^2$}
\opt{C}{$1/a^3$}
\opt{D}{$a^3$}
\end{qoptions}

\qn{19} If $\displaystyle\lim_{x\to\frac{-1}{\sqrt2}}\dfrac{f(x)-2}{2x^2-1}=3\pi$, then $\displaystyle\lim_{x\to\frac{-1}{\sqrt2}}e^{f(x)}$ will be
\begin{qoptions}
\opt{A}{$e^3$}
\opt{B}{e}
\opt{C}{$2e^2$}
\opt{D}{None of these}
\end{qoptions}

\qn{20} The acute angle of intersection between the curves $y=x^2$ and $\begin{cases}x=\dfrac53\cos t\\ y=\dfrac54\sin t\end{cases}$ is
\begin{qoptions}
\opt{A}{$\tan^{-1}\left(\dfrac{41}2\right)$}
\opt{B}{$\tan^{-1}\left(\dfrac2{41}\right)$}
\opt{C}{$\tan^{-1}\left(\dfrac34\right)$}
\opt{D}{None of these}
\end{qoptions}

\sect{Section-II}

\qn{1} If a, 2b, c are distinct numbers in A.P. and line ax + by + c = 0 always passes through a fixed point (p, q), then value of $|p+2q|$ is

\qn{2} If x = log2 and y = log3, then a + bx + cy = [log1 + log(1 + 3) + log(1 + 3 + 5) + ..... + log(1 + 3 + 5 + ..... +19)] $-$ 2[log1 + log2 + log3 + ..... +log7], where a, b and c are positive integers. The value of 2a + 3b + 5c is equal to (where log a = $\log_{10}$a)

\qn{3} Let $\vec p,\vec q,\vec r,\vec s$ be the vectors such that $\vec s=x(\vec p\times\vec q)+y(\vec q\times\vec r)+z(\vec r\times\vec p)$. If $[\vec p\ \vec q\ \vec r]=8$ and $\vec s\cdot\left[\vec p+\vec q+\vec r\right]=16$, then x + y + z is equal to

\qn{4} A is a diagonal matrix of order 6 where all its diagonal entries are positive such that det.(A) = 64. Then minimum value of Tr(A) is

\qn{5} If y = f(x) is continuous and f(3) = 0 and f$'$(3) = 6 then $\displaystyle\lim_{x\to0}\dfrac{f(3+3x)+f(3+5x)}x$ is

\qn{6} The number of roots of equation $z^5=\bar z$ is

\qn{7} Let `A' be the area bounded by the curves $y^2=4k_1x$, $k_1\in\left[\dfrac18,\dfrac14\right]$ and $x^2=4k_2(2-y)$, $k_2\in\left[\dfrac14,1\right]$ and y-axis in the first quadrant. Then the least value of A is

\qn{8} If y(x) is a solution of $(1+x)\dfrac{dy}{dx}=1+xy$ and y(0) = $-1$, then y(1) =

\qn{9} Coefficient of $x^8$ in $(1+x)^{10}+x(1+x)^9+x^2(1+x)^8+\dots+x^9(1+x)+x^{10}$ is

\qn{10} Let $L_1:\dfrac{x-5}3=\dfrac{y-7}{-1}=\dfrac{z+2}1$\\
$L_2:\dfrac{x-3}3=\dfrac{y-3}2=\dfrac{z-6}4$\\
be two lines in space. If $L_3$ is a line with direction ratios $\langle2,7,-5\rangle$ meets $L_1$ and $L_2$ at A and B respectively, then the value of $(AB)^2$ is

\end{multicols}
\end{document}
